\chapter{Introduction}
\label{ch:intro}

\section{Background}
\label{sec:background}

Learning manipulation from demonstrations is limited by the cost of acquiring
robot-executable examples. Teleoperation records robot states and actions
directly, but needs a robot and a control interface during collection. Passive
RGB-D capture instead records a person's unconstrained motion with cameras.
It is cheaper to collect, but it does not itself provide robot actions: hand
motion must be reconstructed, related to objects, and transferred to a robot
with different geometry. Even a plausible joint trajectory need not lift or
place an object when contacts and dynamics are simulated.

This thesis studies that chain as a source of training data. The target
experiment is deliberately narrow: a person transfers a blue $40$~mm cube
from a red support cube to a green target cube; a simulated old-generation
UR10 with a Robotiq 2F-85 gripper attempts the same task. The goal is not to
claim a general-purpose manipulation policy, nor to claim deployment on a
physical robot that is not available for this study.

\section{Problem Statement}
\label{sec:problem}

The central question is whether passive demonstrations provide a measurable
advantage when learning a robot policy in physics simulation. Three distinct
gaps intervene. Perception must recover consistent hand and object motion from
two imperfectly calibrated RGB-D views. Retargeting must turn a human motion
into a reachable robot plan with a meaningful gripper command. Finally, the
plan must survive physical execution: a kinematically valid approach may still
miss a cube, slip during transfer, or fail on release.

Object-relative trajectories offer a way to transfer demonstrations to new
cube placements, but they require a reliable object pose and must be checked
again after transfer. Synthetic variations are therefore candidates for
training, not valid demonstrations merely because a geometric transform can
be computed. A second distinction is between learning with object coordinates
and learning from images. The former isolates the quality of transferred data
and control; the latter also tests perception and the visual simulation gap.

\section{Research Gap}
\label{sec:gap}

Prior work covers parts of this problem. Teleoperation research develops
retargeting objectives, while passive-video systems estimate human and object
motion. MimicGen demonstrates that object-centric subtask transfer plus
simulation filtering can expand a small source dataset \cite{mimicgen}. This
work does not claim those operators as novel. Its question is whether an
accessible, camera-only capture pipeline can produce useful robot-learning
examples for the specified cube-transfer task, and which stage limits the
yield. The comparison must include a controller that uses no ViKi examples;
otherwise success in simulation cannot be attributed to the capture pipeline.

\section{Research Questions and Hypotheses}
\label{sec:rq}

\noindent\textbf{RQ.} \emph{Can passive two-view RGB-D demonstrations processed by
\viki\ improve learning of a simulated UR10 cube-transfer policy, compared
with the same training task without ViKi-derived examples?}

\begin{enumerate}[label=\textbf{RQ\arabic*.}, leftmargin=2.2em]
  \item What fraction of recorded episodes survives hand reconstruction,
        object tracking, retargeting and physical replay, and where are the
        dominant losses?
  \item Does object-aware, physics-validated transfer to new cube placements
        improve held-out task success over equal-size world-frame or hand-only
        training sets?
  \item Under matched observations, control, training budget and evaluation
        placements, do ViKi-derived examples improve policy success or sample
        efficiency relative to learning without those examples?
  \item Can a visual policy using simulated camera frames and robot state
        retain task success without privileged object coordinates at test time?
\end{enumerate}

These are testable hypotheses, not findings. H1 predicts a nonzero yield of
physically successful transferred episodes. H2 predicts that object-aware
examples improve success on unseen, reachable cube placements when example
count is controlled. H3 predicts that prior ViKi examples reduce the number
of simulation interactions needed to reach a given success rate. H4 predicts
that a camera-based policy trained from a successful state-based teacher can
perform the task when evaluated using only deployment-available observations.
Failure at any gate is a result to report, not a reason to relabel a geometric
plan as a physically successful demonstration.

\section{Contributions}
\label{sec:contributions}

The implemented contribution is an offline pipeline from calibrated RGB-D
recordings to hand trajectories, object tracks, and UR10 retarget plans. The
experimental contribution sought here is a reproducible assessment of
processing yield, simulated physical validity, object-relative variation,
and the marginal value of ViKi-derived data for state and visual policies.
Simulation replay, learning outcomes and real-world transfer are not claimed
until measured. The proposed experiment uses no motion-capture ground truth
or physical UR10; geometric consistency metrics must therefore not be
presented as absolute pose accuracy.

\section{Thesis Structure}
\label{sec:structure}

Chapter~\ref{ch:lit} reviews capture, retargeting and demonstration-driven
robot learning. Chapter~\ref{ch:method} defines the implemented pipeline and
the planned simulation study. Chapter~\ref{ch:impl} separates current evidence
from experiments still to run. Chapter~\ref{ch:disc} treats uncertainty and
validity, and Chapter~\ref{ch:concl} states the claims that the eventual
results can and cannot support.
